\ChapterOpener{26}{The Theory of Organizational Design and Culture}{Delegation, Accountability, Decision Rights, and Leadership}{How should a founder arrange people, authority, information, incentives, and shared norms so that an organization can make sound decisions, coordinate specialized work, develop its members, adapt to change, and survive the eventual absence of its founder?}

\begin{quote}
\textit{Governing question: How should a founder arrange people, authority, information, incentives, and shared norms so that an organization can make sound decisions, coordinate specialized work, develop its members, adapt to change, and survive the eventual absence of its founder?}
\end{quote}

\vspace{1.5pt}

\section*{Orientation and Scope}

Hiring identifies potential contributors. Operations designs the flow of work. Organization design determines how contributors exercise authority together; culture determines what they actually do when formal instructions do not settle the question. A tiny founder-led company may initially coordinate through direct conversation. A durable institution cannot rely forever on proximity to the founder, nor can it substitute charts, slogans, and committees for responsible human judgment.

\textbf{Scope and evidence note.} This chapter distinguishes research-supported concepts from founder-level operating prescriptions. It draws on Galbraith's Star Model, Mintzberg's organizing configurations, the economics of incentives, Edmondson's research on team learning, Schein's analysis of culture, decision-rights frameworks, and organizational-change scholarship. Frameworks are diagnostic tools, not proven universal recipes. All Bayes Institute organizational structures, personnel, budgets, and scenarios below are proposed or hypothetical; they should not be mistaken for existing staff, a functioning university, accredited status, or verified business results. Employment obligations and governance requirements depend on jurisdiction and contractual commitments.

\begin{table*}[t]
\centering
\small
\caption{Reading Map for Organizational Design and Culture}
\label{tab:ch26_reading_map}
\begin{tabularx}{\textwidth}{p{4.4cm} X}
\toprule
\textbf{Part} & \textbf{Core Question} \\
\midrule
I. Why organization design matters & How does collaboration create value and coordination cost? \\
II. Architecture and structure & When are functional, product, geographic, matrix, or professional forms appropriate? \\
III. Authority, delegation, and accountability & Who may decide, who must execute, and how are rights bounded? \\
IV. Coordination and management systems & How should plans, meetings, information, and performance reviews work? \\
V. Incentives, behavior, and culture & What causes people to act consistently with, or against, stated values? \\
VI. Leadership and talent development & How can managers build independent capability and successor leaders? \\
VII. Change and organizational evolution & How do structures change without destroying identity or delivery? \\
VIII. Structures across industries & Why must professional, software, service, manufacturing, and research firms differ? \\
IX. The continuing Bayes Institute case & How should a tiny technical school grow toward a credible institution? \\
X. Founder's operating templates & What can be designed and tested immediately? \\
\bottomrule
\end{tabularx}
\end{table*}

\subimport{./}{00-why-organization-design-matters.tex}
\subimport{./}{01-organization-architecture-and-structure.tex}
\subimport{./}{02-authority-delegation-and-accountability.tex}
\subimport{./}{03-coordination-management-and-information.tex}
\subimport{./}{04-incentives-behavior-and-culture.tex}
\subimport{./}{05-leadership-managers-and-talent-development.tex}
\subimport{./}{06-change-adaptation-and-organizational-evolution.tex}
\subimport{./}{07-organization-design-across-industries.tex}
\subimport{./}{08-bayes-institute-organizational-design.tex}
\subimport{./}{09-founders-operating-toolkit.tex}
\subimport{./}{10-advanced-decision-cases-and-foundations.tex}
